\chapter{Validation, Defect Calculus, and Method Promotion}
\label{ch:bf-highdim-validation-defect-promotion}
\label{ch:bf-highdim-candidate-synthesis}

This chapter closes the restarted high-dimensional block.  It synthesizes the
validation and promotion logic shared by the deterministic Gaussian/sparse-grid
lane and the non-Gaussian low-rank retained-object lane.  Its purpose is not to
re-derive the methods, but to consume their exported defects, diagnostics,
approximate-target contracts, and HMC-admissibility consequences.

The previous chapters have already separated:
\begin{itemize}[leftmargin=0.28in]
  \item exact target from approximate target,
  \item carried object from moment or retained-object engine,
  \item frozen branch-defining structure from recomputed numerical values.
\end{itemize}
The role here is therefore downstream and comparative: which diagnostics veto a
method, which diagnostics only explain it, and under what conditions a chapter-
local approximation becomes an industrially credible component rather than a
merely interesting computation.

\section{Validation Architecture And Benchmark Roles}
\label{sec:bf-hd-validation-architecture}

The restarted high-dimensional block uses three ledgers of evidence that must not
be conflated:
\begin{enumerate}[label=(\arabic*)]
  \item \textbf{Engineering correctness.}  Did the declared scalar build, factor,
  and update without violating its own branch rules?
  \item \textbf{Numerical or sampler validity.}  Did finite-difference parity,
  ESS, PSD, rank, support, Jacobian, or divergence diagnostics pass?
  \item \textbf{Scientific interpretation.}  Given the surviving diagnostics,
  what may be concluded about approximation quality, target fidelity, or method
  promotion?
\end{enumerate}
A result in one ledger cannot be silently upgraded into the next.  A green build
is engineering evidence.  A valid same-scalar finite-difference row is numerical
evidence.  Neither by itself is a scientific proof that the approximation is the
right industrial method.

The benchmark roles are therefore tiered:
\begin{itemize}[leftmargin=0.28in]
  \item exact or controlled references test whether the declared scalar behaves
  as intended on small models,
  \item local diagnostics test whether the approximation object stays meaningful,
  \item promotion ladders decide whether a method deserves more expensive
  downstream comparison.
\end{itemize}

\section{Finite-Difference Verification As The First Shared Gate}
\label{sec:bf-hd-validation-fd-first}

The first shared validation gate is finite-difference parity on a fixed branch.
For a base parameter \(\theta_0\) and coordinate direction \(e_i\), define
\begin{equation}
  D_i(h)=
  \frac{\widehat\ell_T(\theta_0+h e_i;B)-\widehat\ell_T(\theta_0-h e_i;B)}{2h},
  \qquad
  G_i=\partial_i\widehat\ell_T(\theta_0;B),
  \label{eq:bf-hd-validation-fd}
\end{equation}
where both perturbed evaluations reuse the same fixed branch ledger and
recompute the parameter-dependent numerical objects through the same declared
equations.  The row is branch-valid only if the branch identity is unchanged on
both sides.

A useful reporting contract is:
\begin{equation}
  e_i(h)=\frac{|D_i(h)-G_i|}{\max\{1,|D_i(h)|,|G_i|\}},
  \label{eq:bf-hd-validation-fd-error}
\end{equation}
with a small ladder such as
\begin{equation}
  h\in\{10^{-1},10^{-2},10^{-3},10^{-4},10^{-5}\}\max\{1,|\theta_i|\}.
  \label{eq:bf-hd-validation-fd-ladder}
\end{equation}
The diagnostic supports the same-scalar claim only when branch-valid rows show a
stable decreasing error window.  Failure may indicate a wrong derivative, a
branch mismatch, or a nonsmooth/ill-conditioned value path.  It is a validation
gate, not a posterior-accuracy theorem.

\section{Why The Synthesis Must Start With Defects}
\label{sec:bf-hd-synthesis-purpose}

The preceding chapters explain high-dimensional nonlinear filtering methods.
An industrial macro-finance program needs a harsher object: a calculus of
failure.  A method is useful for a global-bank or central-bank model only when
the failure mode is explicit, the mitigation has a mathematical contract, and
the cost variable that can kill the method is visible before implementation.
Academic novelty is not the target here.  The target is a defensible
composition of known methods under stated assumptions.

This chapter therefore treats every candidate as a conditional component.  A
Gaussian scaffold is useful if it exposes local curvature and covariance
failure.  Sparse-grid or high-degree cubature is useful if it diagnoses local
nonlinearity without pretending to be a global filter.  Tensor methods are
useful if economic structure keeps ranks stable and preserves probability or
covariance semantics.  Transport maps are useful if they reduce geometry while
leaving a checkable target.  HMC is useful only downstream of a same-scalar
likelihood contract.

The reader-facing synthesis map is this: start with the exact conditional law
and likelihood normalizers from Chapter~\ref{ch:bf-highdim-nonlinear-foundations};
replace them only by named approximations; attach one diagnostic to each object
being approximated; and promote a method only after the previous object exports
the input required by the next.  The chapter's exact target is not a new
algorithm.  It is an auditable decision procedure for composing existing methods
without losing the scalar likelihood, probability-law, support, Jacobian, rank,
or covariance semantics needed by industrial inference.

\paragraph{Running quadratic-observation cell.}
The scalar cell
\[
  x_t=\rho x_{t-1}+\eta_t,\qquad y_t=x_t^2+\epsilon_t
\]
has now served five roles.  Chapter~\ref{ch:bf-highdim-nonlinear-foundations}
used it to show that the exact filtering law can have two lobes.
Chapter~\ref{ch:bf-highdim-gaussian-quadrature} used it to show why Gaussian
projection can preserve moments while losing shape.  Chapter~\ref{ch:bf-highdim-fixed-cloud-sgqf-validation}
used it to turn one fixed sparse-grid cloud into a concrete value-path and
numeric-oracle contract.  Chapter~\ref{ch:bf-highdim-low-rank-density-kr} and
Chapter~\ref{ch:bf-highdim-squared-tt-fixed-branch-likelihoods} used it to
separate retained-object semantics from fixed-branch likelihood construction.
Chapter~\ref{ch:bf-highdim-hmc-research} used it to insist that HMC needs a
scalar likelihood and a same-scalar gradient.  Here the same cell is not
evidence for a macro-finance model.  It is the small mechanism that tells a
mixed numerical reader what the architecture is trying not to lose.

\paragraph{Imports from the preceding chapters.}
\begingroup
\scriptsize
\setlength{\tabcolsep}{2pt}
\begin{longtable}{@{}p{0.14\linewidth}p{0.19\linewidth}p{0.22\linewidth}p{0.17\linewidth}p{0.18\linewidth}@{}}
\toprule
Source chapter & Exported object & Exported defect & Exported diagnostic/veto
& Consumed here as \\
\midrule
\endhead
Ch.~\ref{ch:bf-highdim-nonlinear-foundations} & conditional law, normalizer,
likelihood gradient & wrong target, missing normalizer, value-gradient
mismatch & finite \(Z_t\), mass, positivity, same-scalar check & target and
likelihood gates \\
Ch.~\ref{ch:bf-highdim-gaussian-quadrature} & Gaussian moment pair and
point-rule family & projection bias, exactness without accuracy, point explosion
& innovation, PSD, active dimension, comparator residual & scaffold and local
curvature diagnostic \\
Ch.~\ref{ch:bf-highdim-fixed-cloud-sgqf-validation} & fixed sparse-grid scalar,
numeric oracle, branch record & merge mismatch, signed-weight instability,
branch-invalid finite differences & oracle reproduction, branch-valid FD,
minimum-eigenvalue veto, cloud reproduction & deterministic Gaussian lane
validation \\
Ch.~\ref{ch:bf-highdim-low-rank-density-kr} & retained object, coordinate map,
KR / preconditioning semantics & support failure, Jacobian mismatch,
rank/semantic failure & support/logdet, rank, mass, positivity & retained-object
meaning gate \\
Ch.~\ref{ch:bf-highdim-squared-tt-fixed-branch-likelihoods} & fixed-branch TT/KR
scalar, stored-field contract & wrong saved object, branch-identity loss,
incomplete next-step closure & stored-field audit, retained-marginal closure,
branch-fixed structural check & operational middle gate \\
Ch.~\ref{ch:bf-highdim-hmc-research} & scalar potential, gradient, transformed
target & same-scalar mismatch, divergences, invalid transform & scalar parity,
finite support/Jacobian, energy diagnostics & downstream sampler gate \\
\bottomrule
\end{longtable}
\endgroup

\section{Lane-Specific Validation Ladders Before Promotion}
\label{sec:bf-hd-validation-lane-ladders}

The high-dimensional part now contains two distinct approximation lanes with
different local failure modes, so the validation ladder must stay lane-aware
before the closing promotion comparison begins.

\subsection{FixedSGQF ladder}
\label{sec:bf-hd-validation-sgqf-ladder}

The sparse-grid lane should not be promoted from elegant construction alone.
Before any cross-method comparison, it should pass a ladder that includes:
\begin{enumerate}[leftmargin=0.28in]
  \item deterministic cloud reproduction and duplicate-merge checks,
  \item the one-step numeric oracle from
  Chapter~\ref{ch:bf-highdim-fixed-cloud-sgqf-validation},
  \item branch-valid finite differences on a saved fixed-cloud branch,
  \item signed-weight stability checks through \(P_t^-\), \(S_t\), and carried
  covariance objects,
  \item controlled GHQ / linear-Gaussian comparators where available.
\end{enumerate}
These diagnostics answer whether the implementation is evaluating the declared
fixed-cloud scalar, not whether Gaussian projection has become an exact
high-dimensional posterior method.

\subsection{Retained-object TT/KR ladder}
\label{sec:bf-hd-validation-ttkr-ladder}

The TT/KR lane carries a richer retained object and therefore faces richer local
failure modes.  Before cross-method promotion, it should pass a ladder that
includes:
\begin{enumerate}[leftmargin=0.28in]
  \item coordinate-system and Jacobian-consistency checks,
  \item positivity / nonnegativity and normalization checks for the represented
  object,
  \item rank-growth and conditioning diagnostics,
  \item stored-field and retained-marginal closure checks for the next step,
  \item branch-valid finite differences on the fixed-branch likelihood route.
\end{enumerate}
These diagnostics do not prove that moderate TT rank is always sufficient or that
the adaptive source algorithm is globally smooth.  They only establish whether
the frozen retained-object branch remains a scientifically interpretable scalar
object on the tested route.

\paragraph{Copied-core and branch-status reporting.}
When the derivative claim is part of the result, the TT/KR lane should report not
only a finite-difference error table but also the branch-status fields imported
from Chapter~\ref{ch:bf-highdim-hmc-research}: branch-identity indicators,
copied-core indicators, decreasing-window indicators, and the resulting status
classification.  A small error alone is insufficient if the perturbed values were
computed on different frozen fields or with copied fitted cores.

A useful derivative reporting object is:
\begin{equation}
  \mathcal R_{\rm FD}=
  \{\widehat\ell_T(\theta_0;B),\ G_i,\ (h_i,D_i(h_i),e_i(h_i))_i,
    I_\pm(h_i),K_\pm(h_i),W_i,\text{status}\}.
  \label{eq:bf-hd-validation-fd-report-object}
\end{equation}
Here the report should also retain the scalar values
\(\widehat\ell_T(\theta_0+h_i e_j;B)\) and
\(\widehat\ell_T(\theta_0-h_i e_j;B)\) for each \(h_i\), not merely the final
error summary.  The derivative table is meaningful only when the branch record is
identical in the three columns base / plus / minus.

\paragraph{Benchmark-specific interpretation.}
The benchmark ledger should stay explicit about what each model can and cannot
establish.  A linear-Gaussian benchmark can establish exact-reference agreement
for the declared run contract, but not nonlinear posterior accuracy.  A
stochastic-volatility or predator-prey benchmark can establish long-horizon
nonlinear stability or preconditioning benefit under the declared run contract,
but not universal scalability.  A spatial SIR benchmark can establish whether the
method remains useful at the tested state dimension, horizon, rank ladder, and
basis ladder, but not that all spatial epidemic models admit moderate TT rank.
One-seed demonstrations remain diagnostic unless the result artifact explicitly
justifies a stronger interpretation.

\paragraph{Runtime, memory, and veto interpretation.}
Runtime and memory should be interpreted only after veto diagnostics have passed.
When reported, they should be split into mathematically meaningful components:
forward target evaluation, fitting solves, mass contractions, KR inversion or
sampling work, retained-object storage work, and derivative overhead.  Likewise,
peak memory should separate value-only storage from value-plus-gradient storage.
These are not mere profiling counters; they are the places where the fixed-branch
approximation can become unaffordable or mathematically fragile.

\section{Mathematical Veto Ledger Before Interpretation}
\label{sec:bf-hd-validation-veto-ledger}

The validation report should mark a run as failed before interpreting posterior
plots, trajectory summaries, or method rankings if any mathematical veto fires.
Let
\begin{equation}
  \mathcal F_t=
  \{Z_t,\widehat Z_t,\widehat p_t,C_{t,k},\dot C_{t,k},w_i,D_h,G_{\rm FB}\}
  \label{eq:bf-hd-veto-ledger-objects}
\end{equation}
collect the scalar, tensor, weight, and derivative objects used at time \(t\).
The finite-value veto is
\begin{equation}
  V_{\rm finite}=
  \mathbf{1}\{\exists t,\exists a\in\mathcal F_t:\ a\text{ contains NaN or Inf}\}.
  \label{eq:bf-hd-veto-finite}
\end{equation}
The normalization veto is
\begin{equation}
  V_{\rm norm}=
  \mathbf{1}\{\exists t:\ \widehat Z_t\le0\ \text{or}\ |\int \widehat p_t(z)\dd z-1|>\varepsilon_{\rm norm}\}.
  \label{eq:bf-hd-veto-norm}
\end{equation}
The least-squares conditioning veto is
\begin{equation}
  V_{\rm cond}=
  \mathbf{1}\{\max_{t,k}\kappa(N_{t,k})>\kappa_{\rm max}^{\rm allowed}\}.
  \label{eq:bf-hd-veto-cond}
\end{equation}
The rank-saturation veto is
\begin{equation}
  V_{\rm rank}=
  \mathbf{1}\{\exists t:\ R_t=R_{\rm max}\ \text{and}\ {\rm residual}_t>\varepsilon_{\rm fit}\}.
  \label{eq:bf-hd-veto-rank}
\end{equation}
The KR monotonicity veto is
\begin{equation}
  V_{\rm KR}=
  \mathbf{1}\{\exists j,z_{<j},a<b:\ F_j(b\mid z_{<j})<F_j(a\mid z_{<j})\}.
  \label{eq:bf-hd-veto-kr}
\end{equation}
Finally, the fixed-branch derivative veto is
\begin{equation}
  V_{\rm grad}=
  \mathbf{1}\{B(\beta+h)\ne B(\beta)\ \text{or}\ B(\beta-h)\ne B(\beta)\ \text{or no branch-stable decreasing window holds}\}.
  \label{eq:bf-hd-veto-grad}
\end{equation}
A run should be interpreted only if
\begin{equation}
  V_{\rm finite}+V_{\rm norm}+V_{\rm cond}+V_{\rm rank}+V_{\rm KR}+V_{\rm grad}^{\rm required}=0,
  \label{eq:bf-hd-veto-sum-rule}
\end{equation}
where \(V_{\rm grad}^{\rm required}=V_{\rm grad}\) only when the claim includes
validating the derivative.  The tolerances in these vetoes must be declared
before a run.  A minimal default table is:
\begin{equation}
\begin{array}{c|c|c}
\text{quantity} & \text{default} & \text{meaning}\\
\hline
\varepsilon_{\rm norm} & 10^2\epsilon_{\rm mach}+10\,\varepsilon_{\rm quad} & \text{normalization tolerance}\\
\kappa_{\rm max}^{\rm allowed} & 10^{10}\ \text{in double precision} & \text{least-squares conditioning ceiling}\\
\varepsilon_{\rm fit} & \max(10^{-8},10\,{\rm tol}_{\rm ALS}) & \text{rank-saturation residual ceiling}\\
\varepsilon_{\rm KR} & 10^2\epsilon_{\rm mach}+10\,\varepsilon_{\rm quad} & \text{monotonicity/inversion tolerance}
\end{array}
  \label{eq:bf-hd-veto-defaults}
\end{equation}
These are implementation tolerances, not scientific outcomes.  Changing them
after seeing posterior plots invalidates the benchmark interpretation.

\section[Particle Collapse]{Defect 1: Log-Weight Dispersion Can Become Exponential Without Structure}
\label{sec:bf-hd-defect-particles}

The bootstrap and SIR filters in Chapter~\ref{ch:bf-highdim-particle-transport-tensor}
are honest baselines because their collapse mechanism is visible.  The aim of
this section is not to prove the strongest available particle-collapse theorem.
It is to record a heuristic defect calculus showing why, without locality,
guided proposals, or tempering, the variance scale entering the weights can grow
rapidly with effective dimension.  The sharper asymptotic statements belong to
the cited particle-filter literature; the present derivation is a warning model,
not a complete theorem of exponential collapse.
\begin{equation}
  \operatorname{ESS}=\frac{1}{\sum_{i=1}^N \bar W_i^2}.
  \label{eq:bf-hd-synth-ess}
\end{equation}
Suppose the log-weight decomposes across conditionally independent observation
blocks,
\begin{equation}
  L_i=\sum_{k=1}^{d_{\mathrm{eff}}} \ell_{ik},\qquad
  \E[\ell_{ik}]=\mu,\qquad \Var(\ell_{ik})=\sigma_\ell^2,
  \label{eq:bf-hd-synth-logweight}
\end{equation}
with weak dependence sufficient for
\(\Var(L_i)\asymp d_{\mathrm{eff}}\sigma_\ell^2\).  This is the informal
calculus behind the high-dimensional particle-filter warnings of
\citet{bengtsson2008curse} and \citet{snyder2008obstacles}.

\begin{proposition}[Log-weight variance warning]
\label{prop:bf-hd-particle-collapse-calculus}
If the log-weight variance grows like
\(\Var(L_i)=v_d\asymp d_{\mathrm{eff}}\sigma_\ell^2\), then the associated
importance-weight dispersion scale grows exponentially in \(v_d\) under the
lognormal heuristic below.  In that regime, a particle filter with proposal
independent of the current observation can enter a practically severe collapse
regime unless ensemble size grows very aggressively with that dispersion scale.
\end{proposition}

\begin{proof}[Derivation]
For a lognormal heuristic, \(\E[W]=\exp(\mu_d+v_d/2)\) and
\(\E[W^2]=\exp(2\mu_d+2v_d)\), so
\[
  \frac{\E[W^2]}{\E[W]^2}=\exp(v_d).
\]
The usual importance-sampling variance factor is therefore exponential in
\(v_d\).  Since the denominator of \eqref{eq:bf-hd-synth-ess} is driven by the
largest normalized weights, stable ESS typically requires sample growth that
offsets this exponential dispersion.  The cited papers give sharper asymptotic
statements under stronger assumptions; the monograph uses this calculation only
to justify a heuristic defect warning about log-weight dispersion, not to claim a
full theorem of exponential particle collapse for every proposal/model pair.
\end{proof}

\section[Row-Coverage Collapse]{Defect 2: Fixed-Box Row Designs Collapse Against Concentrated Step Targets}
\label{sec:bf-hd-defect-row-coverage}

The second defect was measured, not anticipated: it was diagnosed on the
squared-TT value ladder of Chapter~\ref{ch:bf-highdim-squared-tt-fixed-branch-likelihoods}
(execution artifacts 2026-08-17 through 2026-08-20) and is recorded here
because its \emph{signature} --- an error plateau that ignores every
resolution knob while in-sample residuals improve --- is exactly the kind of
misleading evidence this chapter's ledgers exist to catch.

\subsection{The mechanism}

The fixed least-squares step fits \(h\) so that \(h^2\) tracks the step
target \(f_t\) on the reference box \(B=[-1,1]^{2n}\), from \(N\) frozen
rows \(z_i\) with weights \(w_i\), while the likelihood increment uses the
\emph{exact} Gram normalizer \(Z_h=\int_B h^2\,d\mu\).  Two structural facts
collide as \(n\) grows.

First, the transition-step target concentrates on the intersection of two
kernel manifolds, \(x_t\approx A x_{t-1}\) (transition) and
\(x_t\approx y_t\) (observation), whose thickness is set by the noise scales
and does \emph{not} grow with the box.  For a uniform-in-\(\mu\) row design
the effective sample size against the target,
\begin{equation}
  \operatorname{ESS}_f
  =\frac{\bigl(\sum_i w_i f_t(z_i)\bigr)^2}{\sum_i w_i^2 f_t(z_i)^2},
  \label{eq:bf-hd-rowcov-ess}
\end{equation}
therefore decays like the volume ratio of a correlated
\(O(\sigma/\text{halfwidth})^{2n}\) slab: geometrically in \(n\).  Measured
on the linear-Gaussian ladder fixture at \(N=8192\) rotated-Sobol rows and
halfwidth \(3\): \(\operatorname{ESS}_f\approx 219\) at \(n=2\) but
\(\approx 11\) at \(n=4\) --- the fit determines
\({\sim}208\) core columns from \({\sim}11\) informative rows.

Second, the exact normalizer converts \emph{unseen} target mass directly
into likelihood bias.  Writing the fitted square as
\(h^2=f_t+e\), the increment carries
\(\log\int_B h^2\,d\mu=\log\bigl(\int_B f_t\,d\mu+\int_B e\,d\mu\bigr)\),
and rows that never visit the mass-bearing manifold produce a fit whose
error term \(e\) removes exactly the mass the rows never saw.  The measured
signature at \(n=4\): a per-step under-count of \(2\)--\(3\) nats at
\emph{every} transition step (the \(t=0\) marginal, whose target is not
concentrated, was exact to \(3\times10^{-3}\)), essentially constant in
rank (\(r=2\ldots8\)), basis degree (\(12\ldots16\)), and row budget
(\(8192\ldots32768\) bought \(21\%\)), while the in-sample weighted residual
\emph{improved} monotonically.  In-sample residuals are computed under the
same defective row measure and are therefore structurally blind to this
defect: they are an explanatory diagnostic only, never a promotion
criterion.

\subsection{Why importance rows make it worse}

The textbook repair --- draw rows from a proposal \(q\) concentrated near
the manifold and reweight by \(1/q\) --- was implemented and \emph{refuted}
at the step level.  With conjugate (exact) Gaussian proposals, defensive
uniform mixtures from \(10\%\) to \(75\%\), and both weighted-\(L^2(\mu)\)
and \(L^2(q\mu)\) fit objectives, the increment error flipped from a
\(-1.7\)-nat under-count to \(+7\) to \(+62\)-nat over-counts, degrading
\emph{monotonically} with basis degree.  The mechanism is a representation
constraint, not a sampling one: once rows resolve the manifold, the
square-root target carries the target's full dynamic range across the fixed
box (relative amplitudes \(e^{10+}\) after the smooth shift), a moderate-rank
polynomial TT cannot represent that spike, its off-manifold oscillation is
unconstrained precisely where rows are sparse, and the exact
\(Z_h=\int h^2\) integrates the oscillation into spurious mass.  The uniform
design had merely \emph{hidden} the spike.  No row design can repair this:
the defect lives in the coordinates, not in the quadrature.

\subsection{The repair: bridged coordinates, not better rows}

The source construction already contains the answer.
\citet{zhao2024ttsequential} (their Section~5) precondition the step: a
bridging density \(\rho_t\) built from already-computed state defines a
Knothe--Rosenblatt change of coordinates \(T_t\) with
\((T_t)_\sharp\rho_t\propto\eta\) for a tensor-product reference \(\eta\),
and the TT fits the \emph{ratio} pushforward
\begin{equation}
  q_{\sharp,t}(u)\;=\;
  \frac{q_t\bigl(T_t^{-1}(u)\bigr)}{\rho_t\bigl(T_t^{-1}(u)\bigr)}\,
  \eta(u),
  \label{eq:bf-hd-rowcov-precond}
\end{equation}
which is \(\eta\) perturbed by the (far less concentrated) ratio
\(q_t/\rho_t\).  Their linear instance (their Section~5.2) takes
\(\rho_t\) Gaussian from estimated moments, so \(T_t\) is a per-step affine
map --- recentring and rescaling the box around where the filter actually
is.  A step-level probe of the affine instance on the \(n=4\) fixture
(block-diagonal maps from filtered moments at \(\kappa=3\) standard
deviations, rows uniform in the new box, Jacobian carried in the target)
moved the step error from \(-1.71\) to \(-0.30\) nats and the coverage
\eqref{eq:bf-hd-rowcov-ess} from \(11\) to \(174\) --- the only lever, among
rows, rank, degree, halfwidth, and importance weighting, that moved the
error at all.  The over-widened \(\kappa=4\) variant degraded again, as the
coverage mechanism predicts.

\subsection{Ledger consequences}

Three rules join the promotion calculus.  (i) A row-design scope must
record a target-coverage diagnostic of the form
\eqref{eq:bf-hd-rowcov-ess}; a fit with collapsed \(\operatorname{ESS}_f\)
is invalid as evidence regardless of its residual.  (ii) In-sample weighted
residuals never gate promotion for scattered-row fits; only same-target
references or resolved-quadrature cross-checks do.  (iii) Fixed-box scopes
do not extrapolate in \(n\): coverage decays geometrically, so each
dimension rung requires either a re-validated coverage diagnostic or the
preconditioned coordinates of \eqref{eq:bf-hd-rowcov-precond}.

The engine-level form of this repair --- per-step frozen triangular maps
propagated through the retained object, with a quadrature correction for
increment mass truncated by the box --- subsequently closed the \(n=2\)
rung outright: \(5.4\times10^{-5}\) nats per step against the declared
\(2.5\times10^{-3}\) bar, with \(r^{*}(2)=6\) confirmed across seeds and
ranks (execution artifacts 2026-08-20 through 2026-08-22).  The same
engine at \(n=4\), however, exposed a third defect --- not in the rows,
and not repaired by the maps --- in the bounded domain itself.
Section~\ref{sec:bf-hd-defect-containment} derives it, and
Section~\ref{sec:bf-hd-refdom} resolves the construction question it
forces.

\section[Chained-Domain Truncation]{Defect 3: Chained Bounded Domains Truncate The Filter Itself}
\label{sec:bf-hd-defect-containment}

The bridged-coordinate repair of
Section~\ref{sec:bf-hd-defect-row-coverage} places every step's fit on
its own moment box.  Doing so exposes a third defect that is invisible at
a single step because it is created by the \emph{recursion}: the retained
polynomial of step \(t-1\) is defined only on the step-\(t-1\) box, so
the step-\(t\) construction is forced to keep its previous-block domain
inside the old one, and the nested domains truncate filter mass at a
geometric rate in the dimension.  This section derives the defect in
proposition--proof form.  Throughout, \(n\) is the state dimension (the
joint fit variable has dimension \(2n\)), \(\Phi\) is the standard normal
distribution function, and ``the engine'' is the adapted squared-TT value
engine with per-step frozen affine maps validated at \(n=2\).

\subsection{Why the domains must nest}

\begin{lemma}[Containment is forced by retention]
\label{lem:bf-hd-containment-nesting}
Let the step-\(t-1\) retained object be a polynomial density
representation defined on the image \(B_{t-1}=m_{t-1}+L_{t-1}[-1,1]^n\)
of its own coordinate map, with no declared extension outside
\(B_{t-1}\).  Then any step-\(t\) fit whose previous-block rows are drawn
in a new box \(B^{p}_{t}=m^{p}_{t}+L^{p}_{t}[-1,1]^n\) is well defined
only if \(B^{p}_{t}\subseteq B_{t-1}\).
\end{lemma}

\begin{proof}
The step-\(t\) target evaluates the retained density of step \(t-1\) at
every previous-block row location \(x_p\in B^{p}_{t}\).  By hypothesis
the retained object assigns no value to \(x_p\notin B_{t-1}\): the
polynomial expression exists there formally but does not represent the
retained density, and the engine's fail-closed assertion rejects such
rows.  A well-defined step therefore requires every point of
\(B^{p}_{t}\) to lie in \(B_{t-1}\).
\end{proof}

The constraint would be harmless if the boxes could remain wide.  They
cannot.  The engine's measured containment behavior (wide-\(\kappa\)
grid, 2026-08-21, two independent measurements) is that the
\emph{effective} previous-block half-width --- the requested \(\kappa\)
times the shrink factor that containment imposes --- is pinned at
\(\kappa^{*}\approx 2.0\)--\(2.2\) standard deviations regardless of the
requested \(\kappa\in[3,6]\): the containment bound involves the ratio of
new-box to old-box widths, both of which scale with the same requested
\(\kappa\), so the request cancels.  We record \(\kappa^{*}\) as a
measured engine constant, not a theorem; the propositions below take it
as an input and are exact given it.

\subsection{The truncated mass is geometric in dimension}

\begin{proposition}[In-box mass of a whitened Gaussian block]
\label{prop:bf-hd-gaussian-box-mass}
Let \(X\sim\calN(m,\Sigma)\) on \(\R^n\) with \(\Sigma=LL^{\top}\)
nonsingular, and let
\(B_\kappa=\{x:\lVert L^{-1}(x-m)\rVert_\infty\le\kappa\}\) be the
moment box with half-width \(\kappa\) in whitened coordinates.  Then
\begin{equation}
  \Pr\{X\in B_\kappa\}=\bigl(2\Phi(\kappa)-1\bigr)^{n}.
  \label{eq:bf-hd-boxmass}
\end{equation}
\end{proposition}

\begin{proof}
Set \(U=L^{-1}(X-m)\).  Then \(U\sim\calN(0,I_n)\), the event
\(\{X\in B_\kappa\}\) is \(\{\lvert U_i\rvert\le\kappa\ \forall i\}\),
the coordinates \(U_i\) are independent standard normals, and
\(\Pr\{\lvert U_i\rvert\le\kappa\}=\Phi(\kappa)-\Phi(-\kappa)
=2\Phi(\kappa)-1\).  The product over \(i=1,\dots,n\) gives
\eqref{eq:bf-hd-boxmass}.
\end{proof}

\begin{proposition}[Geometric growth of the truncated fraction]
\label{prop:bf-hd-truncation-growth}
Under the hypotheses of Proposition~\ref{prop:bf-hd-gaussian-box-mass},
the truncated fraction
\(\varepsilon_{\mathrm{trunc}}(n,\kappa)
 =1-\bigl(2\Phi(\kappa)-1\bigr)^{n}\)
is strictly increasing in \(n\) and satisfies
\(\varepsilon_{\mathrm{trunc}}(n,\kappa)\ge
 1-e^{-2n\overline{\Phi}(\kappa)}\)
with \(\overline{\Phi}(\kappa)=1-\Phi(\kappa)\); in particular it does
not vanish for any fixed \(\kappa\) as \(n\) grows, and at the measured
containment cap it is already material at \(n=4\):
\begin{center}
\begin{tabular}{lccc}
\toprule
 & \(n=2\) & \(n=4\) & \(n=8\)\\
\midrule
\(\kappa=2.0\) & \(0.089\) & \(0.170\) & \(0.311\)\\
\(\kappa=2.2\) & \(0.055\) & \(0.107\) & \(0.202\)\\
\(\kappa=3.0\) & \(0.005\) & \(0.011\) & \(0.021\)\\
\bottomrule
\end{tabular}
\end{center}
\end{proposition}

\begin{proof}
Monotonicity in \(n\) is immediate since
\(0<2\Phi(\kappa)-1<1\).  For the bound, write
\(2\Phi(\kappa)-1=1-2\overline{\Phi}(\kappa)\) and use
\(\log(1-x)\le -x\):
\((1-2\overline{\Phi}(\kappa))^{n}\le e^{-2n\overline{\Phi}(\kappa)}\).
The tabulated values evaluate \eqref{eq:bf-hd-boxmass} directly
(diagnostic script \code{check\_reference\_domain\_selection\_20260823.py}).
\end{proof}

The measured engine telemetry sits \emph{above} this ideal bound, as it
must: the containment shrink is anisotropic and the box is not perfectly
centered.  On the \(n=4\) ladder cells (attempt04, 2026-08-22) the ratio
of increment mass recovered outside the box to mass inside was
\(0.36\)--\(1.85\) across six cells --- in three of six, more step mass
lay outside the box than inside --- against the ideal
\(0.11\)--\(0.17\) of the table.  At \(n=2\) the same telemetry read
\(0.03\)--\(0.11\), which is why the \(n=2\) rung passes with margin.

\subsection{Truncation is a bias floor, not a resolution problem}

The decisive point is that the truncated mass is not an approximation
error that better fitting reduces.  Even a \emph{perfect} fit on the box
retains the wrong law.

\begin{proposition}[Total-variation cost of conditioning]
\label{prop:bf-hd-conditioning-tv}
Let \(P\) be a probability measure and \(B\) an event with
\(P(B)=1-\varepsilon>0\), and let \(P_B=P(\cdot\mid B)\) be the
conditioned law.  Then
\(\operatorname{TV}(P_B,P)
 =\sup_{A}\lvert P_B(A)-P(A)\rvert=\varepsilon.\)
\end{proposition}

\begin{proof}
For any measurable \(A\), decompose \(P(A)=P(A\cap B)+P(A\cap B^{c})\)
and \(P_B(A)=P(A\cap B)/(1-\varepsilon)\).  Hence
\[
  P_B(A)-P(A)
  =P(A\cap B)\Bigl(\tfrac{1}{1-\varepsilon}-1\Bigr)-P(A\cap B^{c})
  =\tfrac{\varepsilon}{1-\varepsilon}P(A\cap B)-P(A\cap B^{c}).
\]
The first term lies in \([0,\varepsilon]\) because
\(P(A\cap B)\le 1-\varepsilon\); the second lies in
\([0,\varepsilon]\).  Their difference therefore lies in
\([-\varepsilon,\varepsilon]\), so the supremum is at most
\(\varepsilon\), and \(A=B^{c}\) attains it:
\(P_B(B^{c})-P(B^{c})=0-\varepsilon=-\varepsilon\).
\end{proof}

\begin{remark}[Zero-fit-error bias floor of the chained-box program]
\label{rem:bf-hd-bias-floor}
Combine the three statements.  Containment
(Lemma~\ref{lem:bf-hd-containment-nesting}) pins the effective
previous-block box at \(\kappa^{*}\approx2\)--\(2.2\) (measured);
at that width a Gaussian block loses the geometric mass fraction of
Proposition~\ref{prop:bf-hd-truncation-growth}; and by
Proposition~\ref{prop:bf-hd-conditioning-tv} the step's retained law then
differs from the true one by at least that fraction \emph{in total
variation even at zero fit error}.  The floor is independent of TT rank,
basis degree, and row budget --- which is exactly the measured attempt04
signature at \(n=4\): errors of \(0.34\)--\(0.47\) nats per step, nearly
flat from rank \(6\) to \(8\), with the truncation-ratio telemetry
identifying the cause.  The one-step statement is exact; the multi-step
compounding (each step conditions anew on a new box) is measured rather
than bounded in closed form here, and it can only worsen the one-step
floor.  A scalar quadrature correction can and does repay the
\emph{increment} mass (it moved the \(n=4\) value error from
\({\sim}2.5\) nats to \(0.34\)--\(0.47\)), but it cannot repair the
\emph{retained shape}: the recursion propagates a minority slice of the
posterior.  This is the precise sense in which the bounded-box program is
insufficient at \(n\ge4\), and it discharges the r\^ole the
truncation-ratio flag was designed for: those cells are invalid as rank
evidence, so no \(r^{*}(4)\) statement of any kind is licensed by them.
\end{remark}

\subsection{Ledger consequences}

Two rules join the calculus.  (i) Any bounded-domain scope whose retained
object chains across steps must record the truncation-ratio diagnostic
(recovered outside-mass over inside-mass); a cell whose ratio exceeds the
declared threshold is invalid as rank or accuracy evidence, whatever its
residuals.  (ii) A domain policy is a structural assumption in the sense
of this chapter's promotion calculus: it must be re-justified at every
dimension rung, because Proposition~\ref{prop:bf-hd-truncation-growth}
makes ``worked at \(n\)'' evidence that expires geometrically in \(n\).

\section[Reference-Domain Selection]{Reference-Domain Selection For The Squared-TT Program}
\label{sec:bf-hd-refdom}

Defect 3 forces a construction question that must not be settled by
preference: on which reference domain, with which basis family and which
reference measure, should the squared-TT recursion live?  This section
states the admissible candidates, defines the properties a construction
must satisfy or optimize, and derives the status of every candidate
against every property.  The conclusion --- recorded here with its
supporting derivations because it redirects the high-dimensional program
--- is that the incumbent bounded construction is \emph{eliminated} by
three independent arguments, and that among the two source-faithful
replacements the Gaussian-reference Hermite construction is selected on
the derived criteria, with the compactified construction retained as the
documented fallback.

\subsection{Candidates and properties}

\begin{definition}[Candidate constructions]
\label{def:bf-hd-refdom-candidates}
All candidates share the per-step frozen affine (triangular)
preconditioner \(T_t(x)=L_t^{-1}(x-m_t)\) built from deterministic
moment hints, the squared representation with defensive term, the exact
Gram-chain normalizer, and the manual-adjoint derivative program.  They
differ only in the triple (reference measure \(\mu\), basis, domain):
\begin{itemize}[leftmargin=0.28in]
  \item[\textbf{C1}] (incumbent) \(\mu=\) uniform on the moment box
    \(B_t\subset\R^{2n}\); Legendre basis; retained domains chained by
    Lemma~\ref{lem:bf-hd-containment-nesting}.
  \item[\textbf{C2}] \(\mu=\calN(0,I_{2n})\) on \(\R^{2n}\) in the
    preconditioned coordinates \(u=T_t(x)\); normalized probabilists'
    Hermite basis; no boxes.
  \item[\textbf{C3}] \(\mu=\) uniform on \([-1,1]^{2n}\) in coordinates
    \(z_i=u_i/\sqrt{s^2+u_i^2}\) (the per-axis algebraic
    compactification used by the reference implementation of
    \citet{zhao2024ttsequential}); Legendre or Lagrange basis on the
    box; no truncation, since the compactification is a diffeomorphism
    of \(\R\) onto \((-1,1)\).
\end{itemize}
\end{definition}

\begin{definition}[Requirements and objectives]
\label{def:bf-hd-refdom-properties}
A construction must satisfy the hard requirements
\begin{itemize}[leftmargin=0.28in]
  \item[\textbf{R1}] \emph{Exact normalizer and analytic adjoint}: the
    one-dimensional mass matrices \(B_{ij}=\int b_i b_j\,d\mu\) are
    closed-form constants, so the Gram-chain normalizer and the manual
    adjoint of the declared finite program exist exactly;
  \item[\textbf{R2}] \emph{Well-defined recursion}: the step-\(t\) row
    law is supported inside the domain of the step-\((t-1)\) retained
    object;
\end{itemize}
and is scored on the objectives
\begin{itemize}[leftmargin=0.28in]
  \item[\textbf{O1}] structural truncation mass
    \(\varepsilon_{\mathrm{trunc}}\) (target mass outside the
    representable domain at zero fit error);
  \item[\textbf{O2}] the zero-fit-error bias floor of the retained law;
  \item[\textbf{O3}] per-axis basis cardinality \(\ell(\varepsilon)\)
    needed for relative \(L^2(\mu)\) accuracy \(\varepsilon\) on the
    program's calibration family (the whitened near-Gaussian step
    targets; the linear-Gaussian model is the exactly-Gaussian member);
  \item[\textbf{O4}] TT rank demand on the same family;
  \item[\textbf{O5}] conditioning of basis evaluations on the effective
    row support, and the \(\ell\)-driven cost channels (design-matrix
    width \(\ell r^2\), mass work \(\ell^2\), compiled-graph size);
  \item[\textbf{O6}] existence of the exact parameter score of the
    declared program with closed-form adjoint nodes;
  \item[\textbf{O7}] the defensive floor performs its declared function:
    the floor density dominates the target
    (\(\sup_x q/\lambda<\infty\)), which is the hypothesis under which
    \citet[eq.~(13) and Lemma~1]{zhao2024ttsequential} prove the bounded
    target-to-approximation ratio.
\end{itemize}
Engineering re-derivation cost is deliberately \emph{not} an objective;
it is a budget input recorded separately, because allowing it to stand in
for a mathematical property is how an inherited construction survives
scrutiny it should not survive.
\end{definition}

\subsection{The reference algebra is closed-form in all three candidates}

\begin{proposition}[Hermite mass identity]
\label{prop:bf-hd-hermite-mass-identity}
Let \(\mathrm{He}_k\) denote the probabilists' Hermite polynomials and
\(\widetilde{\mathrm{He}}_k=\mathrm{He}_k/\sqrt{k!}\).  Under
\(\mu=\calN(0,1)\),
\(\int \widetilde{\mathrm{He}}_j\,\widetilde{\mathrm{He}}_k\,d\mu
 =\delta_{jk}\): the mass matrix is the identity, exactly, at every
degree.
\end{proposition}

\begin{proof}
The generating function of the family is
\(G(z,t)=e^{zt-t^2/2}=\sum_{k\ge0}\mathrm{He}_k(z)\,t^k/k!\).
With \(\varphi\) the standard normal density,
\[
  \int G(z,s)\,G(z,t)\,\varphi(z)\,dz
  =e^{-(s^2+t^2)/2}\int e^{z(s+t)}\varphi(z)\,dz
  =e^{-(s^2+t^2)/2}\,e^{(s+t)^2/2}
  =e^{st},
\]
using the normal moment generating function
\(\int e^{za}\varphi(z)\,dz=e^{a^2/2}\).  Expanding both sides,
\[
  \sum_{j,k}\frac{s^j t^k}{j!\,k!}
  \int \mathrm{He}_j\,\mathrm{He}_k\,d\mu
  =\sum_{m}\frac{(st)^m}{m!},
\]
and matching coefficients of \(s^j t^k\) gives
\(\int \mathrm{He}_j \mathrm{He}_k\,d\mu=\delta_{jk}\,k!\).  Dividing by
\(\sqrt{j!\,k!}\) yields the claim.  (Numerical confirmation: the
degree-\(\le12\) mass matrix computed by 60-node Gauss--Hermite
quadrature matches the identity to \(1.3\times10^{-15}\).)
\end{proof}

\begin{proposition}[C3 leaves the incumbent reference algebra unchanged]
\label{prop:bf-hd-c3-algebra}
Under C3 the reference measure is the uniform probability measure on
\([-1,1]^{2n}\), identical to C1's; every mass matrix, Gram chain,
retention contraction, and adjoint node of the incumbent program is
unchanged.  The compactification enters only through the fit target: with
\(u_i=s z_i/\sqrt{1-z_i^2}\) per axis and Jacobian
\(du_i/dz_i=s(1-z_i^2)^{-3/2}\), the density of a physical target \(q\)
with respect to \(\mu\) in \(z\)-coordinates is
\begin{equation}
  q_\mu(z)\;=\;q\bigl(u(z)\bigr)\prod_{i=1}^{2n}
  s\,(1-z_i^2)^{-3/2}\Big/ 2^{-2n},
  \label{eq:bf-hd-c3-target}
\end{equation}
a row-evaluable factor with no effect on the basis-side algebra.
\end{proposition}

\begin{proof}
The reference algebra of the incumbent program depends only on the pair
(basis, \(\mu\)); C3 keeps both.  Equation \eqref{eq:bf-hd-c3-target} is
the change-of-variables formula for densities under the product
diffeomorphism \(z\mapsto u(z)\), with \(2^{-2n}\) the uniform reference
density.  All \(z\)-dependence introduced by the map sits inside the
target evaluation at rows, which the fit already treats as a black box.
\end{proof}

\begin{remark}[The derivative requirement never forced the box]
\label{rem:bf-hd-no-forcing}
R1 holds for all three candidates ---
Proposition~\ref{prop:bf-hd-hermite-mass-identity} for C2 (where the
identity mass matrix in fact \emph{simplifies} every Gram chain),
Proposition~\ref{prop:bf-hd-c3-algebra} for C3, and the existing
certified program for C1.  R6 below is likewise
candidate-independent.  The program's historical premise that the
analytic-derivative requirement selected the bounded Legendre box is
therefore \emph{wrong relative to that claim}: the only place the
derivative requirement genuinely forces a deviation from
\citet{zhao2024ttsequential} is the preconditioner's moment estimate ---
their Section~5.2 estimates \((\mu_t,\Sigma_t)\) by a particle step,
which is not differentiable, so this program uses frozen deterministic
hints --- and that deviation is identical across C1, C2, and C3.
\end{remark}

\subsection{Domain validity eliminates the incumbent}

\begin{proposition}[Domain scorecard]
\label{prop:bf-hd-domain-scorecard}
(i) C1 fails R2 at the measured containment cap for every unbounded-%
support state model, with truncation mass growing geometrically in \(n\)
(Propositions~\ref{prop:bf-hd-gaussian-box-mass}--%
\ref{prop:bf-hd-truncation-growth}) and a zero-fit-error retained bias of
at least that mass (Proposition~\ref{prop:bf-hd-conditioning-tv}).
(ii) C2 and C3 satisfy R2 with
\(\varepsilon_{\mathrm{trunc}}\equiv0\) and zero structural bias floor:
their domains are all of \(\R^{2n}\) (directly for C2; via the
diffeomorphism for C3), so the previous-step retained object is defined
at every possible row, containment never arises, and at zero fit error
the retained law equals the true one up to the defensive-term mixture
quantified in Proposition~\ref{prop:bf-hd-defensive-mixture-bias}.
\end{proposition}

\begin{proof}
(i) is the content of Defect 3
(Section~\ref{sec:bf-hd-defect-containment}).  For (ii), the retained
object under C2 or C3 is a density on the full space; the step-\(t\) row
law, wherever it is supported, lies in that domain, so
Lemma~\ref{lem:bf-hd-containment-nesting} imposes no constraint.  With
no excluded region, \(\varepsilon_{\mathrm{trunc}}=0\) by definition,
and the zero-fit-error retained density equals
\((q_\mu+\tau Z_h\lambda)/((1+\tau)Z_h)\), whose distance from the truth
is exactly the mixture term treated below.
\end{proof}

\subsection{The defensive floor: bias cost and domination}

The squared representation carries the defensive term of
\citet[eq.~(13), eq.~(16)]{zhao2024ttsequential}: the represented
\(\mu\)-density is \(\bar q=h^2+\tau Z_h\lambda\) with
\(Z_h=\int h^2 d\mu\), \(\lambda\ge0\), \(\int\lambda\,d\mu=1\), so the
increment \(\int\bar q\,d\mu=(1+\tau)Z_h\) remains exact.  The floor
exists to keep the target absolutely continuous with respect to the
representation --- bounded log-ratios, importance weights satisfying a
central limit theorem --- and its two costs are quantified next.

\begin{proposition}[Defensive mixture bias]
\label{prop:bf-hd-defensive-mixture-bias}
At zero fit error (\(h^2=q_\mu\), the \(\mu\)-density of the true step
target with mass \(Z_h\)), the normalized represented law is the mixture
\(\hat p=(1-w)\,p+w\,\lambda\) with \(p=q_\mu/Z_h\) and
\(w=\tau/(1+\tau)\), and
\[
  \operatorname{TV}(\hat p,p)=w\operatorname{TV}(\lambda,p)\le
  \frac{\tau}{1+\tau},
\]
independently of the dimension.  Under the source coupling
\(\tau_t\le\lVert\phi_t-\sqrt{q_t}\rVert_{L^2}^2\)
\citep[Lemma~1 and Proposition~11]{zhao2024ttsequential} the bias is
second order in the fit error, so the zero-fit-error limit of the
retained law is exactly the truth; under this program's current fixed
policy \(\tau=10^{-6}\) it is a per-step constant three orders below the
\(2.5\times10^{-3}\) accuracy bar.
\end{proposition}

\begin{proof}
Normalizing \(\bar q=q_\mu+\tau Z_h\lambda\) by its mass
\((1+\tau)Z_h\) gives
\(\hat p=\frac{1}{1+\tau}p+\frac{\tau}{1+\tau}\lambda\).  Then
\(\hat p-p=w(\lambda-p)\), and integrating absolute values,
\(\operatorname{TV}(\hat p,p)=\tfrac12\int|\hat p-p|
 =w\cdot\tfrac12\int|\lambda-p|=w\operatorname{TV}(\lambda,p)\le w\).
The mixture weight \(w\) contains no dimensional quantity.  The two
\(\tau\) policies substitute directly.
\end{proof}

\begin{proposition}[Bounded domains defeat the floor's own hypothesis]
\label{prop:bf-hd-domination-failure}
Let the step target \(q\) have full support on \(\R^{2n}\) (as for every
Gaussian-family fixture).  Then for any defensive density \(\lambda\)
supported in a bounded box \(B\), \(\sup_x q(x)/\lambda(x)=\infty\), so
the domination hypothesis of \citet[eq.~(13)]{zhao2024ttsequential}
fails and the guarantee
\(\sup_x p/\hat p<\tfrac{\hat z}{\tau z}\sup_x q/\lambda\) is vacuous.
Under C2 and C3 the hypothesis is satisfiable by design: \(\lambda\) may
be any full-support tensor-product density with closed-form mass
(the reference density itself, or a heavier-tailed product chosen to
dominate a declared target tail class), so the floor's tail order is a
free design parameter rather than a property of the construction.
\end{proposition}

\begin{proof}
On \(B^{c}\), \(q>0\) while \(\lambda=0\), so the ratio is unbounded.
Under C2/C3 take \(\lambda\) to be any product density positive on
\(\R^{2n}\) with \(\int\lambda\,d\mu=1\); positivity everywhere makes
\(\sup q/\lambda\) finite for any continuous \(q\) with tails dominated
by \(\lambda\)'s, which is a choice of \(\lambda\), not a constraint of
the domain.
\end{proof}

\begin{remark}[What the floor regularizes, and what it does not]
\label{rem:bf-hd-floor-channels}
The floor splits the tail question into two channels that must not be
conflated.  \emph{Recursion stability}: with domination, the represented
log-density is floored at \(\log(\tau Z_h\lambda)\), so zeros of the
fitted polynomial and un-fitted tail regions produce bounded,
tail-mass-weighted errors at later rows instead of divergent
log-targets or unbounded ratios --- this is exactly the guard's declared
function, available to C2 and C3 equally through the choice of
\(\lambda\), and structurally unavailable to C1 by
Proposition~\ref{prop:bf-hd-domination-failure}.
\emph{Increment accuracy}: the floor's mass is \(\lambda\)-shaped, not
target-shaped; Lemma~1 of \citet{zhao2024ttsequential} prices it as
additional \(L^2\) error, and target mass the fit does not capture still
biases the increment one for one.  The floor therefore does not relax
any approximation-efficiency requirement below; it converts tail
\emph{failures} from catastrophic to graded, which is a safety property,
not an accuracy subsidy.
\end{remark}

\subsection{Approximation efficiency separates the survivors}

\begin{proposition}[Degree collapse at the Gaussian calibration point of C2]
\label{prop:bf-hd-c2-degree-collapse}
Let the step target be \(q(x)=Z\,\calN(x;m,\Sigma)\) on \(\R^{d}\)
(\(d=2n\)), and let the preconditioner use the exact moments:
\(T(x)=L^{-1}(x-m)\), \(LL^{\top}=\Sigma\).  Then the C2 fit target ---
the density of the pushforward \(T_\sharp q\) with respect to
\(\eta=\calN(0,I_d)\) --- is the constant \(Z\), so the exact fit is
\(h\equiv\sqrt{Z}\): per-axis degree \(0\), TT rank \(1\).
\end{proposition}

\begin{proof}
The pushforward density with respect to Lebesgue measure is
\(q(T^{-1}(u))\,\lvert\det L\rvert
 =Z\,\calN(m+Lu;m,\Sigma)\,\lvert\det L\rvert\).
Since
\(\calN(m+Lu;m,\Sigma)
 =(2\pi)^{-d/2}\lvert\det\Sigma\rvert^{-1/2}
  e^{-\frac12 u^{\top}u}\)
and \(\lvert\det L\rvert=\lvert\det\Sigma\rvert^{1/2}\), this equals
\(Z\,(2\pi)^{-d/2}e^{-\lVert u\rVert^2/2}=Z\,\eta(u)\).  Dividing by the
reference density \(\eta\) leaves the constant \(Z\).
\end{proof}

For the linear-Gaussian calibration model the joint step target
\emph{is} Gaussian in \((x_t,x_{t-1})\) --- the product of a Gaussian
retained law, a linear-Gaussian transition, and a linear-Gaussian
observation has a jointly quadratic exponent --- so
Proposition~\ref{prop:bf-hd-c2-degree-collapse} applies exactly under
exact hints, and near-Gaussian targets perturb it continuously: under C2
the basis burden is set by the \emph{deviation} from Gaussianity, which
is the rank-collapse mechanism the source paper demonstrates for its
preconditioned fits (their Figure~2) and the regime its delegated
weighted-approximation theory addresses
\citep{cui2022deep,cui2023selfreinforced}.

\begin{proposition}[No geometric Legendre rate for the compactified Gaussian]
\label{prop:bf-hd-c3-subgeometric}
Under C3 with map scale \(s>0\), the one-axis fit target at the Gaussian
calibration point is
\(f(z)=c\,(1-z^2)^{-3/4}\exp\!\bigl(-a\,\tfrac{z^2}{2(1-z^2)}\bigr)\)
on \([-1,1]\), with \(a=s^2/2\) and \(c>0\) (the square root of the
compactified Gaussian--Jacobian product).  Its Legendre coefficients
\(c_k\) satisfy \(\limsup_{k\to\infty}\lvert c_k\rvert^{1/k}=1\): the
expansion converges at no geometric rate, in contrast to the entire
(supergeometric) fit targets of C1 and the constant fit target of C2 at
the same calibration point.
\end{proposition}

\begin{proof}
Suppose, for contradiction,
\(\limsup_k\lvert c_k\rvert^{1/k}=1/\rho<1\).  By the Laplace integral
representation
\(P_k(z)=\frac1\pi\int_0^\pi\bigl(z+\sqrt{z^2-1}\cos\theta\bigr)^k
 d\theta\),
every Legendre polynomial obeys
\(\lvert P_k(z)\rvert\le\bigl(\lvert z\rvert+\lvert\sqrt{z^2-1}\rvert
\bigr)^k\le\rho'^k\) on the closed Bernstein ellipse
\(E_{\rho'}\) (foci \(\pm1\), parameter \(\rho'\)), for any
\(\rho'<\rho\).  Hence \(\sum_k\lvert c_k\rvert\,\lvert P_k(z)\rvert\)
converges uniformly on \(E_{\rho'}\), and the Legendre series defines a
function \(F\) analytic on the open ellipse \(\operatorname{int}E_\rho\),
which contains the points \(\pm1\) in its interior for every
\(\rho>1\), and which agrees with \(f\) on \((-1,1)\).  Then
\(F^4\) is analytic on \(\operatorname{int}E_\rho\) and agrees on
\((-1,1)\) with
\[
  f(z)^4=c^4\,(1-z^2)^{-3}\exp\!\Bigl(-\tfrac{2a\,z^2}{1-z^2}\Bigr),
\]
which is single-valued and analytic on \(\mathbb{C}\setminus\{\pm1\}\)
(the algebraic factor is a rational function; the exponential factor is
analytic wherever \(z\ne\pm1\)).  By the identity theorem on the
connected open set \(\operatorname{int}E_\rho\setminus\{\pm1\}\), the
two agree there, so \(f^4\) extends analytically to \(z=1\).  But along
\(z=1+\delta e^{i\pi/4}\), \(\delta\downarrow0\), one has
\(1-z^2=-2\delta e^{i\pi/4}+O(\delta^2)\), hence
\[
  -\frac{2a\,z^2}{1-z^2}
  =\frac{a}{\delta}e^{-i\pi/4}+O(1),
  \qquad
  \operatorname{Re}\Bigl(-\tfrac{2az^2}{1-z^2}\Bigr)
  =\frac{a}{\sqrt2\,\delta}+O(1)\to+\infty,
\]
so \(\lvert f^4\rvert\ge C\delta^{-3}e^{a/(\sqrt2\delta)}\to\infty\):
\(f^4\) is unbounded near \(z=1\), contradicting analyticity there.
Therefore no \(\rho>1\) exists.
\end{proof}

\begin{remark}[Quantitative verification and cost consequences]
\label{rem:bf-hd-refdom-quant}
The proposition rules out a geometric rate; the practical degrees were
computed exactly (Gauss--Legendre quadrature of the coefficient
integrals, 2026-08-23 diagnostic script).  For relative \(L^2\) accuracy
\(10^{-3}\) on one axis at the Gaussian calibration point, C3 needs
per-axis cardinality \(\ell=33\) at map scale \(s=1\) (\(63\) at
\(s=0.5\), \(21\) at \(s=2\)); the C1 comparator (entire
\(\sqrt{\text{truncated Gaussian}}\)) needs \(\ell=7\); C2 needs
\(\ell=1\) by Proposition~\ref{prop:bf-hd-c2-degree-collapse}.  The
reference implementation of \citet{zhao2024ttsequential} corroborates
the C3 figure independently: its compactified route runs \(33\) degrees
of freedom per axis.  The \(\ell\) burden propagates through every cost
channel of O5 --- design-matrix width \(\ell r^2\), mass work
\(\ell^2\), and compiled-graph size, the channel whose rank-driven
variant already produced measured compile-time failures at nine axes ---
while C2's recorded O5 cost is the tail growth of normalized Hermite
values, bounded by \({\sim}26\) at \(\lvert z\rvert=4\) for degrees up
to \(16\), a region Gaussian-reference rows essentially never exceed.
\end{remark}

\subsection{Selection, decisions, and what remains open}

\begin{center}
\begin{tabular}{lccc}
\toprule
Property & C1 (box) & C2 (Hermite) & C3 (compactified)\\
\midrule
R1 exact normalizer/adjoint & pass & pass & pass\\
R2 recursion well defined & \textbf{fail}, \(n\ge4\) & pass & pass\\
O1 truncation mass & geometric in \(n\) & \(0\) & \(0\)\\
O2 bias floor & irreducible & \(O(\tau)\), controlled & \(O(\tau)\), controlled\\
O3 \(\ell(10^{-3})\), calibration & \(7\) (via truncation) & \(1\) & \(21\)--\(63\)\\
O5 cost channel & moot & bounded tail factor & \(\ell\)-inflated\\
O7 floor domination & \textbf{fail}, intrinsic & satisfiable & satisfiable\\
\bottomrule
\end{tabular}
\end{center}

C1 is eliminated on three independent grounds: R2
(Lemma~\ref{lem:bf-hd-containment-nesting} with the measured cap), the
geometric truncation identity with its conditioning bias floor
(Propositions~\ref{prop:bf-hd-gaussian-box-mass}--%
\ref{prop:bf-hd-conditioning-tv}), and the intrinsic failure of the
defensive floor's domination hypothesis
(Proposition~\ref{prop:bf-hd-domination-failure}).  Between the
survivors, every derived criterion either ties (R1, R2, O1, O2, O6, O7)
or favors C2 (O3 decisively at and near the calibration point; O5
through the \(\ell\)-driven channels), so the program adopts \textbf{C2},
with C3 recorded as the fallback should the C2 derivation or
implementation encounter a blocking obstacle --- C3 reuses the incumbent
algebra (Proposition~\ref{prop:bf-hd-c3-algebra}) at a permanent
per-axis cost that the reference implementation itself pays.

Three program decisions are recorded with this selection.  First, the
defensive constant moves from the fixed convenience value to the source
coupling \(\tau_t\le\lVert\phi_t-\sqrt{q_t}\rVert_{L^2}^2\) of
\citet[Proposition~11]{zhao2024ttsequential}, which makes the mixture
bias of Proposition~\ref{prop:bf-hd-defensive-mixture-bias} second order;
the fixed value remains for smoke tests only.  Second, the defensive
density \(\lambda\) defaults to the reference density, with the
domination requirement of
Proposition~\ref{prop:bf-hd-domination-failure} to be re-verified
against the declared tail class of each non-Gaussian target, escalating
to a heavier-tailed product \(\lambda\) where domination demands it.
Third, Proposition~\ref{prop:bf-hd-c2-degree-collapse} changes what the
rank ladder measures: under C2 with exact hints the linear-Gaussian
ladder collapses by construction, so it becomes a machine-precision
correctness oracle, and the feasibility question --- the smallest rank
achieving declared accuracy as dimension grows --- moves to targets that
are genuinely non-Gaussian after preconditioning (the stochastic-%
volatility family) or to deliberately degraded hints.  A ladder that
kept the old question under the new construction would pass for the
wrong reason.

What remains open is stated so that it cannot be silently promoted:
the rank demand of C2 on non-Gaussian targets is an empirical question
with no prediction recorded here; the practical Hermite conditioning
constant at \(n\ge8\) and the compiled-kernel behavior of Hermite
evaluations are engineering measurements yet to be made; and the
delegated weighted-approximation theory
\citep{cui2022deep,cui2023selfreinforced} is cited as the source of the
unbounded-domain error analysis and must be read against the C2
derivation before that derivation is treated as reviewed.

\section{Validation Benchmarks And Promotion Logic}
\label{sec:bf-hd-validation-benchmarks}


The validation ladder should combine exact or controlled references, local
object-validity diagnostics, and promotion rules.  Exact or reference-density
metrics such as Hellinger distance, relative \(L^1\) error, moment error,
trajectory RMSE, coverage, and ESS quantiles are meaningful only when the target
and branch contracts have already been named.  Their role is not to replace the
same-scalar contract, but to test whether a declared approximation remains
useful after the engineering and branch-validity gates have passed.

For the restarted high-dimensional block, the benchmark architecture should at
least cover:
\begin{enumerate}[label=(\arabic*), leftmargin=0.28in]
  \item an exact or controlled linear-Gaussian reference where available,
  \item a nonlinear toy case that reveals posterior-shape or retained-object
  limitations,
  \item a branch-valid finite-difference check on the declared approximate scalar,
  \item structural diagnostics specific to the lane: rank, mass, positivity,
  PSD, support, Jacobian, or ESS as appropriate.
\end{enumerate}
These ladders are not optional appendices.  They are what stops a chapter-local
approximation from being promoted on speed or elegance alone.

The benchmark backbone should be named explicitly so the reader does not have to
reconstruct it from isolated examples:
\begin{enumerate}[label=(V\arabic*), leftmargin=0.34in]
  \item exact-reference linear Gaussian benchmark,
  \item long-horizon stochastic-volatility benchmark,
  \item spatial SIR benchmark for high-dimensional partial-observation filtering,
  \item predator-prey benchmark for nonlinear preconditioning,
  \item derivative validation table whenever a fixed-branch gradient is part of
  the claim.
\end{enumerate}
A useful compact benchmark table should pair the resource ledger with the
model-specific accuracy ledger and, when derivatives are claimed, with a separate
finite-difference table containing \(D_h\), \(G_{\rm FB}\), \(|D_h-G_{\rm FB}|\),
and the branch-status fields for each \(h\).  The derivative table is meaningful
only when the branch record is identical in the three columns
\(\theta-h\), \(\theta\), and \(\theta+h\).

\paragraph{What each benchmark can and cannot establish.}
The linear-Gaussian benchmark can establish exact-reference agreement for the
declared run contract, but it cannot establish nonlinear posterior accuracy.  The
stochastic-volatility benchmark can establish long-horizon nonlinear stability and
credible agreement with references or synthetic truth under the declared run
contract, but it cannot by itself establish high-dimensional scalability.  The
spatial SIR benchmark can establish whether the method remains useful at the
tested latent dimension, horizon, rank ladder, and basis ladder, but it cannot
prove that all spatial graphs or epidemic models admit moderate TT rank.  The
predator-prey benchmark can establish that nonlinear preconditioning yields a
measurable benefit on the tested nonlinear dynamical system, but it cannot prove
that such preconditioning is always worth its cost.  One-seed demonstrations are
diagnostic unless the result artifact explicitly justifies stronger
interpretation.

\section{Defect Calculus And Promotion Rules}
\label{sec:bf-hd-promotion-logic}

The chapter closes with a simple rule: runtime, memory, ESS per gradient, rank,
and point count can compare methods only after the target, semantic, and
numerical veto diagnostics have passed.  The industrial promotion logic is
therefore conditional:
\begin{equation}
  V \text{ passes } \Longrightarrow \text{compare performance};
  \qquad
  V^c \Longrightarrow \text{veto, not ranking.}
  \label{eq:bf-hd-validation-veto-logic}
\end{equation}
This chapter does not conclude that high-dimensional nonlinear filtering is
solved for BayesFilter, that any method is NAWM-ready, that HMC converges, or
that any lane should become a production default.  It concludes only that a
candidate method must pass the declared validation and defect calculus before
speed, elegance, or novelty are allowed to matter.
